% SYNC-ID: discussion
\section{Discussion}\label{sec:discussion}
% CONTENT_DEFERRED_TO_WRITING_PHASE
\textit{(THIS IS PLACEHOLDER TEXT, PLZ FULFILL IT WHEN WRITING.)}
% SYNC-ID: discussion-placeholder-1
Discussion will be written after claims, evidence, limitations, and epistemic strength have been reviewed. The current paragraph only reserves the expected reading order.
% SYNC-ID: discussion-placeholder-2
Observed patterns will later be connected to limitations without adding a new claim in the discussion.
% SYNC-ID: discussion-placeholder-3
This block reserves a calibrated interpretation and leaves its content to the approved evidence record.
% SYNC-ID: discussion-placeholder-4
Open statements will be identified explicitly after the review pass.
% SYNC-ID: discussion-placeholder-5
The scope comparison will be short, paired, and limited to what the study supports.
% SYNC-ID: discussion-placeholder-6
The final wording will keep causal language calibrated rather than implying unsupported mechanisms.
% SYNC-ID: discussion-placeholder-7
Each interpretation will receive an evidence pointer before this placeholder is accepted.
% SYNC-ID: discussion-placeholder-8
The final discussion paragraph is reserved for a measured transition to validity and conclusion. It should make the boundary between observation, interpretation, and open question easy to audit.
% SYNC-ID: discussion-placeholder-9
This fixture also reserves a short paragraph for noting what the current record cannot establish. The wording will be chosen together with the corresponding limitation, not improvised during translation.
% SYNC-ID: discussion-placeholder-10
The discussion layout includes a little additional space for a reader to follow the move from an observation to a qualified interpretation. A finished version should identify the source of that move and retain any uncertainty that matters. This generic sentence is only a typesetting fixture and does not assert that an observation, interpretation, or open question currently exists.
